\documentclass[11pt]{article}

\usepackage[margin=1in]{geometry}
\usepackage{amsmath,amssymb,booktabs,tabularx}

\title{Long-variance unit-mean normalization}
\author{}
\date{}

\begin{document}
\maketitle

\noindent\textbf{Status.} Draft\\
\textbf{Plan item.} Long-variance unit-mean normalization\\
\textbf{Depends on.} None\\
\textbf{Feeds.} P12, first-moment guard, affine schedule transform, terminal pricing measure

\section{Purpose}

Derive the mean-reversion level that makes the physical long-variance state
have stationary first moment one.
This fixes the scale of \(Z_t\), so persistent
variance memory changes return-clock exposure without introducing an
unidentified unconditional variance multiplier.

\section{Setup}

Work under the physical measure \(P\).
The long-variance state follows
\[
dZ_t
=
\kappa_z(\bar z-Z_t)\,dt
+\sigma_z\sqrt{Z_t}\,dW_t^Z
+\chi_+\,dN_t^+
+\chi_-\,dN_t^-.
\]

The event feedback terms are uncompensated positive jumps. Let
\(\lambda_{\mathrm{stat}}^+\) and \(\lambda_{\mathrm{stat}}^-\) denote the
stationary buy and sell event rates under the mean-normalized seasonal baseline
exposure and the expected news activity loading used for the interval family
being normalized.

\section{Notation}

\begin{tabularx}{\textwidth}{lXX}
\toprule
Symbol & Meaning & Units or domain\\
\midrule
\(Z_t\) & Long-variance state & Positive\\
\(\kappa_z\) & Mean-reversion rate & Positive inverse time\\
\(\bar z\) & Mean-reversion level & Positive\\
\(\sigma_z\) & Volatility of the long-variance diffusion & Non-negative\\
\(\chi_+\) & Buy-event long-variance jump loading & Non-negative\\
\(\chi_-\) & Sell-event long-variance jump loading & Non-negative\\
\(\lambda_{\mathrm{stat}}^+\) & Stationary buy event rate & Non-negative inverse time\\
\(\lambda_{\mathrm{stat}}^-\) & Stationary sell event rate & Non-negative inverse time\\
\(b_Z\) & Stationary event-feedback drift into \(Z_t\) & Inverse time\\
\bottomrule
\end{tabularx}

\section{Assumptions used}

\begin{tabularx}{\textwidth}{XXX}
\toprule
Assumption & Where it is used & Why it is needed\\
\midrule
The Hawkes activity system has finite stationary side event rates.
& Taking expectations of \(dN_t^\pm\).
& The first-moment equation needs finite \(\lambda_{\mathrm{stat}}^\pm\).\\
Event loadings \(\chi_\pm\) are deterministic non-negative constants.
& Pulling jump sizes outside event-count expectations.
& Positive event feedback must enter as a linear first-moment drift.\\
The Brownian innovation has zero conditional mean.
& Taking expectations of \(\sigma_z\sqrt{Z_t}\,dW_t^Z\).
& The diffusion term should not move the first moment.\\
The normalization target is \(\mathbb E^P[Z_t]=1\) in stationarity.
& Solving for \(\bar z\).
& This identifies the scale of the long-variance factor.\\
\bottomrule
\end{tabularx}

\section{Result}

\noindent\textbf{Identity.} Define the stationary event-feedback drift
\[
b_Z
=
\chi_+\lambda_{\mathrm{stat}}^+
+\chi_-\lambda_{\mathrm{stat}}^-.
\]

The mean-reversion level that gives \(Z_t\) unit stationary first moment is
\[
\bar z
=
1-\frac{b_Z}{\kappa_z}
=
1-
\frac{
\chi_+\lambda_{\mathrm{stat}}^+
+\chi_-\lambda_{\mathrm{stat}}^-
}{\kappa_z}.
\]

\section{Derivation}

Let \(m(t)=\mathbb E^P[Z_t]\).
Taking first moments in the long-variance state
equation gives
\[
\frac{d}{dt}m(t)
=
\kappa_z\bar z
-\kappa_z m(t)
+\chi_+\lambda_t^+
+\chi_-\lambda_t^-.
\]

In stationarity, \(\lambda_t^\pm=\lambda_{\mathrm{stat}}^\pm\) at the
first-moment level and \(m(t)=m_\infty\), so
\[
0
=
\kappa_z\bar z
-\kappa_z m_\infty
+\chi_+\lambda_{\mathrm{stat}}^+
+\chi_-\lambda_{\mathrm{stat}}^-.
\]

Solving for the stationary first moment gives
\[
m_\infty
=
\bar z
+
\frac{
\chi_+\lambda_{\mathrm{stat}}^+
+\chi_-\lambda_{\mathrm{stat}}^-
}{\kappa_z}
=
\bar z+\frac{b_Z}{\kappa_z}.
\]

This should not be a second unconditional variance scale,
so the normalization sets \(m_\infty=1\), so that \(Z_t\) is a dimensionless
relative variance multiplier.
Therefore
\[
1
=
\bar z+\frac{b_Z}{\kappa_z},
\]
and hence
\[
\bar z
=
1-\frac{b_Z}{\kappa_z}.
\]

The diffusion parameter \(\sigma_z\) does not appear in the first-moment
normalization because the Brownian term has zero conditional mean.

\section{Domain and validity}

The identity is valid when
\[
\kappa_z>0,\qquad
\sigma_z\ge 0,\qquad
\chi_+\ge 0,\qquad
\chi_-\ge 0,
\]
the stationary side event rates are finite and non-negative, and the derived
level satisfies
\[
\bar z>0.
\]

The strict positivity condition is equivalent to the first-moment guard derived
separately:
\[
\kappa_z>
\chi_+\lambda_{\mathrm{stat}}^+
+\chi_-\lambda_{\mathrm{stat}}^-.
\]

This is not a full statement that all higher moments exist, only a unit first-moment normalization.

\section{Consequences}

\begin{itemize}
    \item The long-variance factor has an identified unit scale under the physical measure.
    \item Positive event feedback does not raise the unconditional mean of \(Z_t\), because \(\bar z\) is lowered by exactly the stationary event-feedback contribution.
    \item The baseline return-clock exposure can use \(Z_t\) as a persistent relative variance multiplier around deterministic baseline opportunity.
    \item Any admissible calibration or pricing contract must carry enough information to derive the stationary side event rates used in \(b_Z\).
\end{itemize}

\section{Checks}

\begin{tabularx}{\textwidth}{XXX}
\toprule
Check & Expected result & Evidence location\\
\midrule
Exact first-moment substitution
& \(\kappa_z(\bar z-1)+b_Z=0\) after substituting the derived \(\bar z\).
& \texttt{check\_long\_variance\_first\_moment.py}\\
Kernel stationary-level contract
& \texttt{stationary\_mean\_reversion\_level} returns \(1-b_Z/\kappa_z\).
& Unit test \texttt{stationary\_mean\_reversion\_level\_offsets\_event\_feedback\_drift}\\
Exogenous-source contract
& Stationary rates include exogenous Hawkes target source rates before deriving \(\bar z\).
& Unit test \texttt{stationary\_mean\_reversion\_level\_includes\_exogenous\_hawkes\_source\_rates}\\
\bottomrule
\end{tabularx}

\end{document}
